% !PRINTED 104
\weberpage{104}
%<<CONT
angedeutet sind, und deren Zahl man ebenso wie oben
\[
= \frac{3^{\mu - 2} - 1}{2}
\]
findet. Wenn also $m$ durch 8 theilbar ist, so ist die Gesammtzahl der Gruppen $\mathfrak{A}$
\weq{5}{\frac{3^{\mu - 1} - 1}{2} + \frac{3^{\mu - 2} - 1}{2} = 2 . 3^{\mu - 2} - 1.}
Die kleinsten Werthe, die $m$ in diesem Falle annehmen kann, sind
\[
m = 8, 12, 15, 20, 21, 24, 28, 33, 35, 39, 40, 44 \ldots
\]
Die kleinste Zahl, die zu mehr als einer solchen Gruppe Anlass giebt, ist $m = 24$, für die man $\mu = 3$ hat, und die also nach (5) fünf verschiedene Gruppen giebt. Man erhält nämlich folgende zulässige und von einander verschiedene Combinationen der $e_{h,k}$: % [sic: Weber again prints e_{h,k}; the displayed combinations and defining congruences use e_{k,h}.]
\[
\begin{array}{ccccc}
1\ 0\ 1 & 1\ 1\ 0 & 1\ 1\ 0 & 1\ 0\ 0 & 1\ 0\ 0\\
0\ 1\ 0, & 0\ 0\ 1, & 0\ 1\ 1, & 0\ 1\ 1, & 0\ 0\ 1,
\end{array}
\]
was zu folgenden Bestimmungen über die Indices führt:
\[
\begin{aligned}
& \alpha_1 \equiv \alpha_3,\quad \alpha_2 \equiv 0,\\
& \alpha_1 \equiv \alpha_2,\quad \alpha_3 \equiv 0,\\
& \alpha_1 \equiv \alpha_2 \equiv \alpha_3,\qquad (\mathrm{mod}\ 2)\\
& \alpha_1 \equiv 0,\quad \alpha_2 \equiv \alpha_3\\
& \alpha_1 \equiv 0,\quad \alpha_3 \equiv 0,
\end{aligned}
\]
und wenn man also die Zahlen $a$ durch die Congruenzen
\[
a \equiv (-1)^{\alpha_2} 5^{\alpha_1} (\mathrm{mod}\ 8), \quad a \equiv 2^{\alpha_3} (\mathrm{mod}\ 3)
\]
bestimmt, so ergeben sich folgende fünf zweigliedrige Gruppen:
\[
\begin{aligned}
\mathfrak{A}_1 &= 1,\ 5\\
\mathfrak{A}_2 &= 1,\ 19\\
\mathfrak{A}_3 &= 1,\ 11\\
\mathfrak{A}_4 &= 1,\ 23\\
\mathfrak{A}_5 &= 1,\ 7.
\end{aligned}
\]
Um die Perioden $\eta$ zu berechnen, können wir ähnlich verfahren, wie im vorigen Paragraphen. Wir bezeichnen mit $\nu_1, \nu_2, \ldots \nu_\mu$ die Indices einer Zahl $n$ und setzen
\[
\begin{aligned}
e_{1,1} \nu_1 + e_{1,2} \nu_2 + \cdots + e_{1,\mu} \nu_\mu &= \sigma_1,\\
e_{2,1} \nu_1 + e_{2,2} \nu_2 + \cdots + e_{2,\mu} \nu_\mu &= \sigma_2,
\end{aligned}
\]
%CONT>>
